% 恒星光谱（综述）
% license CCBYSA3
% type Wiki

本文根据 CC-BY-SA 协议转载翻译自维基百科\href{https://en.wikipedia.org/wiki/Stellar_classification}{相关文章}。

在天文学中，恒星分类是根据恒星的光谱特性对恒星进行的分类。来自恒星的电磁辐射通过棱镜或衍射光栅被分解成光谱，呈现出类似彩虹的连续颜色，其间夹杂着光谱线。每一条光谱线都对应一种特定的化学元素或分子，而谱线的强度则反映该元素的丰度。不同光谱线强度的变化主要由恒星光球层的温度决定，尽管在某些情况下也确实存在真实的元素丰度差异。恒星的光谱型是一种简短的代码，主要概括了其电离状态，从而为光球层温度提供一个客观的度量。

目前，大多数恒星采用摩根–基南（Morgan–Keenan，MK）分类系统进行分类，该系统使用字母 O、B、A、F、G、K、M，按从最热（O 型）到最冷（M 型）的顺序排列。每一个字母类别又进一步用一个数字细分，其中 0 表示最热，9 表示最冷（例如 A8、A9、F0、F1 就构成了一个由热到冷的序列）。此外，该序列还扩展出 W、S、C 三个类别，用于那些不符合经典体系的恒星。一些恒星遗迹或质量明显偏离恒星范围的天体也被赋予了字母标记：D 表示白矮星，L、T、Y 表示棕矮星（以及部分系外行星）。

在 MK 系统中，还会在光谱型后加上一个用罗马数字表示的光度级。光度级依据恒星光谱中某些吸收线的宽度来判定，这些线宽随恒星大气密度变化，因此可用于区分巨星与矮星。其中，光度级 0 或 Ia+ 表示特超巨星，I 表示超巨星，II 表示亮巨星，III 表示普通巨星，IV 表示亚巨星，V 表示主序星，sd（或 VI）表示亚矮星，而 D（或 VII）表示白矮星。太阳的完整光谱分类为 G2V，这表明太阳是一颗主序星，其表面温度约为 5800 K。
\subsection{传统颜色描述}
传统的恒星颜色描述只考虑恒星光谱的峰值位置。然而实际上，恒星会在整个电磁波谱范围内辐射。由于所有光谱颜色叠加后呈现为白色，人眼实际看到的恒星颜色要比传统颜色描述所暗示的浅得多。这种“明亮度”的特性说明，将恒星简单地赋予某种光谱颜色可能具有误导性。在排除暗光条件下的颜色对比效应后，在一般观测条件中，并不存在真正的绿色、青色、靛色或紫色恒星。像太阳这样的“黄色”矮星实际上呈现为白色；“红矮星”呈现为较深的黄色或橙色；而“棕矮星”并不会真的呈现棕色，而是假想中在近距离观察时可能显得暗红色，或呈灰色/黑色。
\subsection{现代分类}
现代恒星分类体系被称为摩根–基南（Morgan–Keenan，MK）分类系统。在该系统中，每一颗恒星都会被赋予一个光谱型（源自较早的哈佛光谱分类体系，该体系本身并不包含光度级$^{[1]}$），并再配以一个用罗马数字表示的光度级（如下所述），二者共同构成恒星的完整光谱类型。

其他现代恒星分类体系（例如 UBV 系统）则基于颜色指数，即通过测量两个或多个波段星等之间的差值来进行分类$^{[2]}$。这些数值通常以诸如 U−V、B−V这样的形式给出，表示通过两种标准滤光片（例如紫外、蓝光和可见光）所测得的颜色差异。
\subsubsection{哈佛光谱分类}
哈佛体系是一种一维的恒星分类方案，由天文学家安妮·江普·坎农（Annie Jump Cannon）提出，她对德雷珀（Draper）早期按字母顺序排列的分类体系进行了重新排序与简化（见“历史”部分）。在该体系中，恒星根据其光谱特征被划分为若干组，每一组用一个字母表示，并可辅以数字细分。

主序星的表面温度大致分布在 2000–50 000 K 之间，而演化程度更高的恒星——尤其是新近形成的白矮星——其表面温度则可超过 100 000 K$^{[3]}$。从物理意义上看，这些光谱类别反映的是恒星大气的温度，并通常按由热到冷的顺序排列。
\begin{figure}[ht]
\centering
\includegraphics[width=14.25cm]{./figures/1615d7b930f17d79.png}
\caption{} \label{fig_Stcl_1}
\end{figure}
用于记忆光谱型字母由热到冷排列顺序的传统助记口诀是“Oh, Be A Fine Guy/Girl: Kiss Me!”（“哦，做个好小伙/好姑娘：吻我吧！”）$^{[12]}$。
尽管在天文学课程和相关组织举办的活动中，人们提出过许多替代性的助记口诀，但这一传统口诀至今仍然是最为流行的$^{[13][14]}$。

光谱类型 O 到 M（以及后文讨论的其他更为专门的类型）都会再用阿拉伯数字 0–9进行细分，其中 0 表示该类型中最热的恒星。例如，A0 表示 A 型中最热的恒星，而 A9 则表示最冷的 A 型恒星。允许使用小数；例如，恒星 Mu Normae 的光谱型被定为 O9.7 $^{[15]}$。太阳的光谱型为 G2$^{[16]}$。

尽管哈佛光谱分类实际上反映的是恒星的表面温度或光球温度（更准确地说是其有效温度），但这一点在该体系建立之初并未被完全理解。不过，到第一张赫罗图（Hertzsprung–Russell diagram）于 1914 年前后提出时，人们已经普遍怀疑这一对应关系是成立的$^{[17]}$。在 1920 年代，印度物理学家梅格纳德·萨哈（Meghnad Saha）在物理化学中分子解离理论的基础上，发展出了电离理论并将其推广到原子的电离问题。他首先将该理论应用于太阳色球层，随后又应用于恒星光谱$^{[18]}$。

随后，哈佛天文学家塞西莉亚·佩恩（Cecilia Payne）证明O–B–A–F–G–K–M这一光谱序列本质上就是一个温度序列$^{[19]}$。由于该分类序列的建立早于人们对其“温度本质”的认识，将一条光谱归入某一具体亚型（如 B3 或 A7）在很大程度上依赖于对恒星光谱中吸收特征强度的（带有一定主观性的）估计。因此，这些亚型并不是在数学意义上等间隔划分的。
\subsubsection{摩根–基南（Morgan–Keenan）分类}
耶基斯光谱分类，也称为 MK 分类，或 Morgan–Keenan（亦称 MKK，Morgan–Keenan–Kellman）系统$^{[20][21]}$，是一种由威廉·威尔逊·摩根（William Wilson Morgan）、菲利普·C·基南（Philip C. Keenan）和伊迪丝·凯尔曼（Edith Kellman）于 1943 年在耶基斯天文台提出的恒星光谱分类体系$^{[22]}$。这是一个二维分类体系（温度与光度），其依据是对恒星温度和表面重力敏感的光谱线；而表面重力又与恒星的光度密切相关（相比之下，哈佛分类仅基于表面温度）。在 1953 年对标准星表和分类标准进行若干修订之后，该体系被正式命名为摩根–基南分类（MK）$^{[23]}$，并一直沿用至今。

表面重力更高、密度更大的恒星，其光谱线会表现出更显著的压力展宽。由于巨星的半径远大于质量相近的矮星，其表面重力（以及由此产生的压力）要低得多。因此，光谱中的差异可以被解释为光度效应，从而仅通过分析光谱本身就能够确定恒星的光度级。

目前区分出了若干不同的光度级，具体列于下表中$^{[24]}$。
\begin{figure}[ht]
\centering
\includegraphics[width=14.25cm]{./figures/73151fcf04a39c9e.png}
\caption{} \label{fig_Stcl_2}
\end{figure}